\emph{The local ACE benchmark.}
Fix \(n\).  The published bound in \(\text{lem:jms-ace-theorem-five-four}\) is valid for every positive value of its cumulant-separation parameter.  Apply that result with the numerical substitution \(\delta=\delta_n\).  Under this substitution the model class \(\mathcal P_{\mathrm{ACE},n}^{\mathrm{JMS}}\) becomes exactly \(\mathcal C_n(\delta_n)\): every restriction in \(\text{def:jms-ace-class}\) is unchanged except that
\[
 |\kappa_{r+1}(\eta)|\geq\delta
 \quad\hbox{is replaced by}\quad
 |\kappa_{r+1}(\eta)|\geq\delta_n.
\]
The quantities \(a_{1,n}\), \(b_{1,n}\), and \(a_2\) do not involve \(\delta\), whereas the quantity \(b_{2,n}\) becomes
\[
 \widetilde b_{2,n}
 =200\min\{1,C_{\theta}\}\delta_n
 \max\!\left\{
   \varepsilon_{1,n},\varepsilon_{2,n},
   (\gamma n)^{-1/2}(\psi_{\xi}+C_{\theta}\psi_{\eta})
 \right\}^{-1}.
\]
Thus the local eligibility condition in the statement is precisely the eligibility hypothesis of \(\text{lem:jms-ace-theorem-five-four}\) after the same substitution.  The assumed positivity of \(\varepsilon_{1,n}\) and \(\varepsilon_{2,n}\), together with the explicit requirement that every logarithm and denominator be defined, supplies all boundary conditions required by that cited result.  Its conclusion is therefore
\[
\begin{aligned}
 &\sup_{P\in\mathcal C_n(\delta_n)}
 Q_{P,1-\gamma}
 \!\left(\left|\widehat\theta_{\mathrm{ACE},r,n}-\theta_0(P)\right|\right)\\
 &\quad\le
 C_{\gamma}r!16^r\delta_n^{-1}
 \Bigl[
   \varepsilon_{1,n}^r\varepsilon_{2,n}
   +C_{\theta}\varepsilon_{1,n}^{r+1}\\
 &\hspace{42mm}
   +64(C_g+\psi_{\eta})^r
    \{r^2(C_g+\psi_{\eta})+\psi_{\xi}+C_{\theta}\psi_{\eta}\}
    (\gamma n)^{-1/2}
 \Bigr],
\end{aligned}
\]
which is the claimed value \(U_{\mathrm{ACE},n}^{\mathrm{loc}}\).  A separately established DML bound and this ACE bound are statements about their respective procedures under their respective hypotheses.  Taking the smaller displayed number is consequently an oracle comparison of two valid upper benchmarks; no data-dependent rule choosing between the procedures, and no lower bound for the local class, follows from these two inequalities alone.

\emph{The Gaussian--Rademacher path.}
Now let \(k=4\), \(r=3\), and \(0<a\leq1\).  Independence of \(G_{\circ}\) and \(S\) gives, for every \(z\in\mathbb C\),
\[
 M_a(z)
 =E[e^{z\sqrt{1-a^2}G_{\circ}}]E[e^{zaS}]
 =\exp\!\left(\frac{1-a^2}{2}z^2\right)\cosh(az).
\]
Cumulants add for independent sums and scale homogeneously.  The Gaussian fourth cumulant is zero, while
\[
 \kappa_4(S)=E[S^4]-3E[S^2]^2=1-3=-2.
\]
Hence \(\kappa_4(\eta_a)=-2a^4\).  The characteristic function is
\[
 \varphi_a(t)=M_a(it)
 =\exp\!\left(-\frac{1-a^2}{2}t^2\right)\cos(at).
\]
The exponential factor is strictly positive on the real line, so its first positive zero is \(t_a=\pi/(2a)\).  Therefore, for every \(d\in\mathbb R\),
\[
 E[\sin\{t_a(\eta_a+d)\}]
 =\operatorname{Im}\{e^{it_ad}\varphi_a(t_a)\}=0.
\]

We next establish relevance.  Differentiating the displayed characteristic function at its zero gives
\[
 \varphi_a'(t_a)
 =-a\exp\!\left\{-\frac{(1-a^2)\pi^2}{8a^2}\right\}
 =-A_a.
\]
For a fixed real \(d\), differentiation under the expectation is justified because \(\eta_a\) has finite moments of every order, and hence
\[
\begin{aligned}
 E[(\eta_a+d)\sin\{t_a(\eta_a+d)\}]
 &=-\operatorname{Re}\left[
   \left.\frac{d}{dt}\{e^{itd}\varphi_a(t)\}\right|_{t=t_a}
 \right]\\
 &=-\operatorname{Re}\{e^{it_ad}\varphi_a'(t_a)\}
 =A_a\cos(t_ad).
\end{aligned}
\]
By \(\eta_a\perp X\), the same identity may be conditioned on \(D_n\) and integrated, yielding
\[
 m_a:=E_P[Z_n\sin(t_aZ_n)]
 =A_aE_P[\cos(t_aD_n)].
\]
Because \(s\geq r+1=4\), Lyapunov's inequality and \(\text{ass:treatment-code-radius}\) imply
\[
 E_P|D_n|\leq\|D_n\|_{L^s(P_X)}leq\varepsilon_{1,n}.
\]
The elementary inequality \(1-\cos u\leq |u|\) now gives
\[
 E_P[\cos(t_aD_n)]
 \geq1-t_aE_P|D_n|
 \geq1-t_a\varepsilon_{1,n}
 \geq\frac12.
\]
Thus \(m_a\geq A_a/2>0\), which also states explicitly the positivity needed by the ratio.

It remains to prove the uniform risk bound.  Put
\[
 W=Z_n\sin(t_aZ_n),
 \qquad
 V=Y\sin(t_aZ_n),
 \qquad
 R=(Y-\theta_0Z_n)\sin(t_aZ_n).
\]
The residual representation gives \(Y-\theta_0Z_n=b_n(X)+\xi\).  Conditional shift annihilation shows
\[
 E_P[b_n(X)\sin(t_aZ_n)]
 =E_P\!\left[b_n(X)E_P\{\sin(t_aZ_n)\mid X\}\right]=0.
\]
Moreover, \(\sin(t_aZ_n)\) is measurable with respect to \((X,T)\), so \(\text{ass:outcome-mean-independence}\) gives
\[
 E_P[\xi\sin(t_aZ_n)]
 =E_P\!\left[\sin(t_aZ_n)E_P(\xi\mid X,T)\right]=0.
\]
Consequently \(E_P[R]=0\) and \(E_P[V]=\theta_0m_a\); this is the identifying moment, distinct from the estimator defined in the statement.

The support conditions provide uniform second-moment constants.  Indeed, clipping and \(\text{ass:g-range}\) imply \(|D_n|\leq2C_g\).  Since \(E[\eta_a^2]=1\) and \(\eta_a\perp D_n\),
\[
 E_P[W^2]\leq E_P[Z_n^2]leq 1+4C_g^2=:K_W.
\]
Also \(|b_n(X)|\leq C_q+2C_{\theta}C_g\).  The conditional sub-Gaussian assumption implies \(E_P[\xi^2]\leq c_{\psi}\psi_{\xi}^2\) for a universal constant \(c_{\psi}\) associated with the \(\psi_2\)-norm convention.  Therefore
\[
 E_P[R^2]
 \leq 2(C_q+2C_{\theta}C_g)^2+2c_{\psi}\psi_{\xi}^2
 =:K_R.
\]
Both \(K_W\) and \(K_R\) are finite and independent of \(a\), \(n\), and the particular path law.

Write \(\overline W=\mathbb P_nW\), \(\overline V=\mathbb P_nV\), and let \(\mathcal E=\{\overline W\geq A_a/4\}\).  On \(\mathcal E\), projection onto \([-C_{\theta},C_{\theta}]\) is a contraction and \(\theta_0\) belongs to that interval, whence
\[
 |\widehat\theta_{a,n}-\theta_0|
 \leq\left|\frac{\overline V}{\overline W}-\theta_0\right|
 =\frac{|\mathbb P_nR|}{\overline W}
 \leq\frac{4}{A_a}|\mathbb P_nR|.
\]
By independent and identically distributed sampling and \(E_P[R]=0\),
\[
 E_P[(\mathbb P_nR)^2]=\frac{\operatorname{Var}_P(R)}{n}leq\frac{K_R}{n}.
\]
On the complementary event the estimator is zero, so its squared error is at most \(C_{\theta}^2\).  Since \(m_a\geq A_a/2\), Chebyshev's inequality yields
\[
 P(\mathcal E^c)
 \leq P\!\left(|\overline W-m_a|\geq\frac{A_a}{4}\right)
 \leq\frac{16K_W}{nA_a^2}.
\]
Combining the last three displays gives
\[
 E_P[(\widehat\theta_{a,n}-\theta_0)^2]
 \leq\frac{16(K_R+C_{\theta}^2K_W)}{nA_a^2}.
\]
Independently, both the estimator and the target lie in \([-C_{\theta},C_{\theta}]\), so the same risk is at most \(4C_{\theta}^2\).  Taking the smaller bound, absorbing fixed numerical factors, and using
\[
 A_a^{-2}
 =a^{-2}\exp\!\left\{\frac{(1-a^2)\pi^2}{4a^2}\right\}
\]
proves the asserted uniform risk inequality with a constant depending only on \(C_{\theta},C_g,C_q,\psi_{\xi}\).

Finally, \(\Delta_a=2a^4\) implies
\[
 a=2^{-1/4}\Delta_a^{1/4},
 \qquad
 a^{-2}=\sqrt2\,\Delta_a^{-1/2}.
\]
Thus \(t_a\varepsilon_{1,n}\leq1/2\) is equivalent to
\[
 \varepsilon_{1,n}\leq\frac{a}{\pi}
 =\frac{\Delta_a^{1/4}}{2^{1/4}\pi}.
\]
Furthermore,
\[
 a^{-2}\exp\!\left\{\frac{(1-a^2)\pi^2}{4a^2}\right\}
 =\sqrt2e^{-\pi^2/4}\Delta_a^{-1/2}
 \exp\!\left\{\frac{\pi^2\sqrt2}{4}\Delta_a^{-1/2}\right\},
\]
and the fixed leading factor is absorbed into \(C\).  This proves the final display.  The argument is confined to the explicit Gaussian--Rademacher path and supplies no lower bound over the larger class defined only by fourth-cumulant separation, so it establishes exactly the stated construction-specific benchmark and no full minimax phase diagram.